\subsection{Définition structurelle et fonction de la constante de Boltzmann}
\label{sec:ii-definicion-boltzmann}

\begin{definition}[Transducteur thermique de l'énergie par décision]
Une fois fixée une section thermique positive
\(\Theta_{{\rm clk},a}\in\mathcal L_\Theta\), le transducteur
de Boltzmann est l'application linéaire positive
\(k_B:\mathcal L_\Theta\to\mathcal L_E\) qui réalise
\begin{equation}
 k_B(\Theta_{{\rm clk},a})
 =\frac{\epsilon_a}{\log3},
 \qquad \epsilon_a=\frac{h_a}{108t_0}.
 \label{eq:ii-definicion-boltzmann-transductor}
\end{equation}
L'action et la durée déterminent \(\epsilon_a\) ; la décision
tritique uniforme détermine sa normalisation informationnelle.
\end{definition}

L'unicité, démontrée à la section~\ref{sec:ii-kb-aut-desarrollo},
résulte de la prescription de l'image d'une section non nulle d'une droite.
Les équations \eqref{eq:ii-kb-aut-coordenada} et
\eqref{eq:ii-kb-aut-orientaciones} précisent, respectivement, sa
coordonnée dans des bases déclarées et la compatibilité d'un même
transducteur avec les deux orientations. La détermination du produit
\(k_B\Theta_{{\rm clk},a}\) précède l'individualisation :
celle-ci conserve comme données explicites la section thermique, les
bases et, lorsque la sélection interne est revendiquée, son lecteur constructeur.
La pression de \eqref{eq:ii-kb-aut-presion}, sous les hypothèses
qui y sont énumérées, normalise le paramètre énergétique
\(\beta=(k_B\Theta)^{-1}\).

La fonction physique de \(k_B\) consiste à exprimer dimensionnellement
l’information d’une distribution : \(S=k_Bs\). Évaluer cette
entropie à une température produit l’énergie \(\Theta S\).
Un triplet uniforme apporte \(\log3\) nats ; une distribution de
continuations apporte son entropie pondérée ; une réduction quantique
apporte l’entropie de la bipartition spécifiée. Le transducteur
agit sur ces contenus en conservant leurs domaines. La mémoire
généalogique détermine quelle information retient chaque lecture ; le
bilan thermodynamique déclare quelle opération l’échange ou
l’efface. Cette séparation explique la relation entre réversibilité,
intrication et énergie de décision.
